\cleardoublepage{}
\begin{center}
    \bfseries \zihao{3} 摘~要
\end{center}

三维高斯泼溅近年来显著推动了高保真新视角合成的发展，能够在保持较高渲染质量的同时实现实时渲染。然而，现有 3DGS 方法主要面向透视相机和 SfM 稀疏点云初始化设计，直接应用于室内移动采集的鱼眼图像时仍面临强畸变投影和长序列标定等问题。本文提出一套基于鱼眼相机与 LiDAR 几何先验的室内三维高斯泼溅重建管线，在保持标准 3DGS 高效训练和渲染框架的同时提升鱼眼场景重建质量并降低长序列标定时间开销。首先，本文利用 LiDAR-IMU 轨迹和 LiDAR-相机预标定外参为鱼眼 SfM 提供位姿先验，减少长轨迹鱼眼序列对逐帧增量式注册的依赖，并降低全局标定的时间开销。其次，本文自动估计鱼眼有效视场角与虚拟视图分辨率，将每张鱼眼图像拆分为八张共享光心的小视场角虚拟透视图，从而避免直接在强非线性鱼眼投影下进行高斯近似。最后，本文引入基于 K-d tree 的 LiDAR 点云下采样与尺度感知初始化，利用雷达点云的局部协方差为高斯提供表面感知的初始尺度和旋转，从而提升后续重建质量。本文在 ScanNet++ 和 Zip-NeRF 等公开鱼眼场景上进行实验，结果表明该管线相比传统去畸变训练、通用畸变相机 3DGS 方法和直接射线空间渲染方法均取得更好的新视角合成质量；同时，快速 SfM 实验表明，本文迭代式标定管线能够在提高标定效率的同时取得与传统 SfM 管线相当的标定效果，消融实验进一步验证了虚拟透视图、位姿二次精化和 LiDAR 初始化对重建质量的贡献。

\bigskip
\noindent \textbf{关键词}：三维高斯泼溅；鱼眼相机；新视角合成；运动结构恢复；激光雷达

\cleardoublepage{}
\begin{center}
    \bfseries \zihao{3} Abstract
\end{center}

3D Gaussian Splatting has recently made significant progress in high-fidelity novel view synthesis, achieving real-time rendering while maintaining high rendering quality.
However, existing 3DGS methods are mainly designed for perspective cameras and SfM-based sparse point initialization. When directly applied to fisheye images captured by indoor mobile systems, they still face challenges such as strong projection distortion and long-sequence calibration. This thesis proposes an indoor 3D Gaussian Splatting reconstruction pipeline based on fisheye cameras and LiDAR geometric priors. The pipeline improves fisheye scene reconstruction quality and reduces the calibration time cost of long image sequences while preserving the efficient training and rendering framework of standard 3DGS. First, it uses LiDAR-IMU trajectories and pre-calibrated LiDAR-camera extrinsics as pose priors for fisheye SfM, reducing the dependence of long fisheye sequences on frame-by-frame incremental registration and lowering the time cost of global calibration. Second, it automatically estimates the effective fisheye field of view and virtual-view resolution, and decomposes each fisheye image into eight small-FoV virtual perspective views with a shared optical center, thereby avoiding direct Gaussian approximation under highly nonlinear fisheye projection. Finally, it introduces K-d tree based LiDAR point-cloud downsampling and scale-aware initialization, which uses local covariance in the LiDAR point cloud to provide surface-aware initial scales and rotations for Gaussians and thereby improves subsequent reconstruction quality. Experiments on public fisheye scenes from ScanNet++ and Zip-NeRF show that the proposed pipeline achieves better novel-view synthesis quality than traditional undistortion-based training, general distorted-camera 3DGS methods, and direct ray-space rendering methods. Meanwhile, fast SfM experiments show that the proposed iterative calibration pipeline improves calibration efficiency while achieving calibration quality comparable to the traditional SfM pipeline, and ablation studies further validate the contributions of virtual perspective views, secondary pose refinement, and LiDAR initialization to reconstruction quality.

\bigskip
\noindent \textbf{Keywords}: 3D Gaussian Splatting; fisheye camera; novel view synthesis; structure from motion; LiDAR
